\newpage
\thispagestyle{empty}
\begin{spacing}{1.2}
\chapter{Introduction}

This chapter provides the background context for the GramaGIS project, outlining the administrative challenges faced by rural Panchayats and the motivation behind developing an intelligent geospatial management system. It establishes the need for integrating GIS and Natural Language Processing technologies to modernize local governance infrastructure.

\section{Background}
Rural administrative units, such as Grama Panchayats, are the foundational layers of governance and local development. As these regions undergo rapid demographic and infrastructural shifts, the need for data-driven planning becomes critical. Managing a Panchayat involves overseeing a complex web of spatial and non-spatial assets, including ward boundaries, educational institutions, healthcare facilities, public amenities, transportation infrastructure, and citizen-reported issues.

Historically, rural administration has relied on physical registers, static paper maps, and fragmented digital spreadsheets. While these methods served traditional needs, they lack the interoperability and real-time analytical capabilities required for modern ``Smart Village'' initiatives. Digital transformation in this sector is often hindered by the technical gap between complex Geographic Information Systems (GIS) and the non-technical officials or citizens who need to use them.

A Geographic Information System (GIS)-based approach, integrated with Natural Language Processing (NLP), provides a structured solution to these challenges. By centralizing disparate datasets into a unified geospatial environment, infrastructure data can be visualized, queried, and managed through an interactive interface.

GramaGIS is developed as an implementation of an Intelligent Rural Infrastructure Management Framework. It leverages a robust spatial stack consisting of PostgreSQL/PostGIS and GeoServer and integrates a Large Language Model (LLM) via the Gemini API to bridge the technical divide. This allows users to interact with complex spatial datasets using natural language, establishing a scalable foundation for transparent and efficient local governance.

GramaGIS was built as a full-stack web GIS portal specifically developed for Vazhathope Grama Panchayath. The system digitizes village-level infrastructure including ward boundaries, road networks, healthcare facilities, educational institutions, commercial services, and government offices and stores them in a centralized PostGIS-enabled database. A GeoServer instance publishes these datasets as OGC-compliant WMS, WFS, and WFS-T services, which are consumed by an interactive Leaflet.js-based map interface. The Node.js and Express.js backend handles authentication, feedback management, GeoServer proxy operations for protected workflows, and the Gemini-assisted natural language query flow. Together, these components deliver a platform where Panchayat officials can manage, query, review feedback, and export spatial data entirely through a web browser, without requiring traditional desktop GIS expertise.

\section{Report Organization}
The remainder of this report is organized as follows. Chapter~2 defines the problem statement and lists the specific objectives that guided the development of GramaGIS. Chapter~3 presents a literature survey covering GIS in rural governance, natural language interfaces for spatial databases, and user-centered design in Web-GIS systems. Chapter~4 describes the system architecture, tracing the evolution from a direct GeoServer connection to the final hybrid Express.js-assisted design. Chapter~5 details the core system features, including the interactive map, intelligent query system, administrative data management module, and feedback workflow. Chapter~6 covers the implementation, explaining the data preparation, database configuration, backend development, and frontend integration steps. Chapter~7 presents a use case analysis illustrating how different actors interact with the platform. Chapter~8 reports the results and achievements mapped against each project objective, supported by system screenshots. Chapter~9 outlines future directions and potential enhancements to the platform. Chapter~10 concludes the report with a summary of contributions and the broader impact of GramaGIS on rural governance.

\end{spacing}
